\newpage
\begin{center}
  \textbf{\large 2. МЕТОДЫ ПЕРЕХВАТА СЕТЕВОГО ТРАФИКА МОБИЛЬНЫХ ПРИЛОЖЕНИЙ ЧЕРЕЗ ПРОКСИ-СЕРВЕР}
\end{center}
\refstepcounter{chapter}
\addcontentsline{toc}{chapter}{2. МЕТОДЫ ПЕРЕХВАТА СЕТЕВОГО ТРАФИКА МОБИЛЬНЫХ ПРИЛОЖЕНИЙ ЧЕРЕЗ ПРОКСИ-СЕРВЕР}


\section{Введение}
% 11 строк

При наличии способов обратной разработки для определения строения мобильного приложения с закрытым исходным кодом и методов анализа серверной части приложения при помощи инструментов автоматического сканирования, обеспечивающих отдельных анализ двух подсистем мобильного приложения, метод перехвата сетевого трафика для определения защищённости мобильного приложения представляется промежуточным решением.

Сетевой слой коммуникации между клиентской и серверной стороной был выбран для дальнейшего исследования по ряду причин, которые делают его наиболее информативным и при этом требующим наименьшее количество затрат времени, ресурсов и проникновения в систему.

Во-первых, в отличие от серверных сервисов и клиентских архитектур, коммуникационный слой более переносимый, стандартизованный, и свободный в смысле порядка и условий передачи пакетов данных от одного участника сети к другому. Это означает, что при смене архитектуры или типа мобильного приложения метод проверки его на защещённость изменится лишь незначительно, что позволит переиспользовать наработки при анализе каждого следущего приложения.

Во-вторых, располагаясь между клиентской и серверной частью приложения, сетевой слой содержит информацию о поведении обоих систем. Анализ сетевого слоя позволяет реализовывать адаптеры между закрытыми от исследования частями приложения и открытыми, постепенно открывая всё больше информации. Это позволяет систематизировать и упростить подход к аудиту мобильного приложения, а также предоставить удобные для проверки промежуточные стадии перевода системы на открытую архитектуру.

В-третьих, использование прокси-сервера посередине предоставляет до этого ограниченную возможность модификации существующей системы, добавления в неё более защищённых модулей вместо уязвимых или полный перенос системы на взаимодействие с разработанным открыто клиентом или сервером. Прокси-сервер также позволяет встроить в систему способ выбора между альтернативными реализациями клиентских и серверных частей, демократизируя доступ в открытый интернет.

Таким образом, именно метод анализа с использованием прокси-сервера посередине между клиентской и серверной частью мобильного приложения приводит к сочетанию лучших возможностей анализа клиентской и серверной части по-отдельности.

\section{Анализ сетевого трафика}

\subsection{Определение адреса серверной части мобильного приложения и поверхностный анализ защитных механизмов}
% 15 строк

Чтобы оправдать применимость метода попробуем получить минимальную ценную информацию в виде ip адреса сервера, обслуживающего мобильное приложение.

Для всех методов анализа мобильного приложения на защищённость кроме автоматического сканирования сервера требуется иметь локальную копию клиентского приложения. Методы получения доступа к недоступному в России клиенту находятся за рамками данной работы, поэтому для проверки работы метода было в первую очередь использовано одно из самых популярных приложений по заказу еды из ресторанов, доступное для установки обе самые распространённые мобильные платформы iOS и Android. Для воспроизводимости эксперимента на персональных компьютерах с популярными операционными системами, была выбрана версия приложения для Android, так как методы эмуляции Android устройств доступны для ОС Windows, Linux и MacOS. Имея собранное Android приложение и виртуальную машину под управлением ОС Android, можно удостовериться в работоспособности приложения в локальном окружении.

Для определения адреса серверной части мобильного приложения достаточно проанализировать весь сетевой трафик в системе, внутри которой запущено мобильное приложение. Пакеты, перехваченные локальным инспектором содержат в себе адрес сервера в качестве адреса получателя, как в стандарте.

После получения ip адреса серверной части приложения можно проверить при помощи обратного поиска DNS домен, по которому доступен API серверной части мобильного приложения из интернета. В случае если обратный поиск DNS найдёт домен, можно составить суждение о примерной структуре приложения. Даже если обратный поиск DNS не находит домен, если мобильное приложение действительно использует домен для обращения к серверу, а не прямой ip адрес, этот домен можно найти в метаинформации запросов, перехваченных к интересующему ip адресу или в локальной таблице DNS записей.

После определения адреса сервера можно попробовать посетить его через веб браузер и найти в заголовках запросов информацию о версии балансировщика нагрузки или обратного прокси-сервера, используемого приложением. Здесь же можно сделать дополнительную проверку на доступность стандартных сервисов или автоматическое сканирование портов, но как уже обсуждалось, это ни к чему не приведёт в случае нестандартных протоколов и архитектурных решений.

\subsection{Получение доступа к зашифрованному трафику}
% 16 строк

Для получения доступа к шифрованному трафику требуется подменить системный TLS/SSL сертификат. Большинство гражданских мобильных приложений полагаются на системные сертификаты и стандартное HTTPS шифрование, частое исключение составляют банковские приложения, проверяющие дополнительно эмитента сертификата, что ограничивает применимость метода к банковским приложениям. Этот метод защиты от атаки типа человек посередине известен как "закрепление сертификата" - при подмене сертификата сервер мобильного приложения отклоняет соединение и не позволяет передать какие-либо данные в обе стороны.

После подмены сертификата мобильное приложение начинает использовать локальный сертификат и становится возможным расшифровать все сетевые пакеты, доступ к которым был получен в предыдущей подсекции.

Доступ к этим данным позволяет проводить аудит безопасного использования пользовательских данных. Используя принцип наименее полного набора данных можно распознать попытки приложения получить доступ к данным, не относящимся к выполнению основной функции, а направленных на сбор пользовательских данных без согласия.

\subsection{Сбор секретных и персонально идентифицирующих данных}

Одним из приложений полученного доступа к перехвату сетевого трафика является сбор секретных, персональных, персонально идентифицирующих или критичных для работы приложения данных. Помимо данных, передающихся при загрузке, старте и работе приложения, интересно наблюдать за данными, передающимися в фоновой работе приложения или при сворачивании приложения, это позволяет провести полноценный аудит безопасности.

Часть данных в приложении загружается с сервера из закрытой базы данных, а доступ к сетевому трафику позволяет перехватывать эти данные и понимать внутреннее устройство приложения. Здесь же можно проверять приложение на наличие недокументированных возможностей и попытаться получить данные о других пользователях системы. Если кто угодно может прочитать персональные данные любого пользователя, за этой прозрачностью стоит риск утечки ценных данных.

Отдельное преимущество метода анализа защищённости мобильных приложений при помощи атаки типа человек посередине заключается в незаметности метода для администрации мобильного приложения. Единственное отличие модифицированного трафика заключается в особом сертификате, но если серверная часть мобильного приложения не проверяет сертификат, то вероятность, что она сохраняет сертификаты для аудита пренебрежительно мала. Таким образом, можно целиком маскировать полезный трафик под органический, немодифицированный трафик обычных пользователей и не испытывать ограничений на скорость и количество соединений, которые появились бы при автоматическом сканировании сетевых интерфейсов.

\subsection{Альтернативные конфигурации перехвата сетевого трафика}

В целях обобщения метода на другие архитектуры мобильных приложений, были также проведены дополнительные эксперименты.

Приложение эксклюзивного бизнес клуба с условием вступления - наличие годовой прибыли от 100 млн. рублей было проверено на возможность утечки персональных данных. Это приложение работает на ОС iOS и было проверено на физическом устройстве. Из-за изменения конфигурации метод не изменился и получилось также перехватывать целиком весь трафик с устройства. Такое применение даёт одну из возможностей для получения доступа к недоступному в России приложению через доступ к устройству с уже установленным приложением.

С такой же целью проверки безопасности персональных данных был разработан особый прокси-сервер для популярной игры Minecraft. Особенность этого клиента заключается в нестандартном с точки зрения интернет технологий методе шифрования трафика. Однако, не смотря на закрытый код серверной части приложения, код для алгоритма шифрования является открытым, что позволяет разработать прокси-сервер, с его поддержкой. Реализация перехвата трафика для такого экзотического протокола подтверждает переносимость метода и его утилитарность.

\section{Результаты}
% 101 строки
% 12 строк

Проведённый эксперимент, сравнивающий подход к анализу защищённости мобильных приложений при помощи перехвата сетевых данных прокси-сервером с обратной разработкой клиентской части и автоматическим сканированием серверной части приложений, показывает какие преимущества и недостатки присутствуют у полученного метода.

Создание слепка сетевых пакетов позволяет полностью воссоздать работу клиентского приложения с точки зрения интерфейса, используемого при общении по сети с серверной частью. Таким образом, как и в случае полной обратной разработки приложения, можно реализовать аналог приложения с закрытым исходным кодом и формально верифицировать точность. Сохраняется возможность создания альтернативных клиентов, что демократизирует доступ к приложениям.

При обратной разработке мобильного приложения требуется тонкая настройка среды разработки, индивидуальная подготовка к исследованию каждой поддерживаемой платформы, а также доступ к собранному приложению и его исполняемому коду. Тогда как при перехвате сетевого трафика процесс аудита выглядит одинаково для каждой платформы с заменой лишь платформо-зависимых элементов, отвечающих за эмуляцию рассматриваемой архитектуры. Более того, анализ возможен даже при отсутствии доступа к исполняемому коду напрямую при наличии устройства, на котором мобильное приложение уже установлено и запущено. Это преимущество ускоряет процесс анализа защищённости и делает данный метод более предпочтительным в ситуациях с ограниченными вычислительными, человеческими и временными ресурсами, и как следствие повышает масштабируемость метода на применение к нескольким интересующим приложениям одновременно.

Как и в случае автоматического сканирования сервера мобильного приложения, перехват коммуникации, идущей к серверу, позволяет делать обоснованные предположения о его работе, узнавать версии установленного серверного программного обеспечения и тестировать систему на защищённость от известных уязвимостей в стандартных протоколах, конфигурациях и сервисах. При использовании разработанного метода сохраняется уверенность в качестве проверки серверной части системы мобильного приложения.

Перехват сетевых данных также собирает всю передающуюся конфиденциальную, персональную, секретную или персонально идетенфицируемую информацию и позволяет более точно проверять возможность утечки данных аналогично с автоматическим сканированием на уязвимости и SQL инъекции. Дополнительно использование прокси-сервера позволяет прогнозировать небезопасное хранение и доступ к данным в случае соблюдения всех требуемых условий на безопасность архитектуры, но наличия ошибок, сбоев или неверного истолкования данных с серверной стороны приложения. Таким образом, помимо полного покрытия возможных уязвимостей, перехват сетевых пакетов предоставляет больше полезной информации к интерпретации чем автоматическое сканирование системы на наличие уязвимостей и SQL инъекций.

Как результат, в сравнении с двумя отдельными методами, рассматривающими лишь часть системы обеспечивающей работу мобильного приложения, выбранный метод позволяет судить о поведении системы в целом, сохраняет преимущества в детекции проблем безопасности и клиентской, и серверной частей приложения, и при этом проще в реализации и масштабировании, чем комбинация автоматического сканирования сетевых интерфейсов сервера и ручного процесса обратной разработки клиентского исполняемого кода. Переносимость метода между разными архитектурами и типами приложений подтверждает универсальную применимость и практическую значимость.

\section{Заключение главы}
% 6 строк

После целостного рассмотрения метода анализа защищённости мобильных приложений при помощи перехвата сетевого трафика промежуточным прокси-сервером, можно сделать вывод о применимости и универсальности полученной стратегии аудита. Возможность полностью рассмотреть систему и сделать обоснованные предположения о внутренней структуре приложения с закрытым исходным кодом, найти потенциальные проблемы сделанные на этапе построения архитектуры или реализации программного обеспечения и возникающие при этом утечки персональных данных.

Полученный метод сохраняет все преимущества конкурирующих методов, устраняет некоторые их недостатки и предоставляет уникальное преимущество в виде переносимости между разными архитектурами. При этом отдельное внимание при выборе метода анализа было уделено сложности самого метода, так как от сложности воспроизведения зависит воспроизводимость метода в промышленных мастштабах.

еосле тестирования на разных платформах, разных архитектурах поддерживаемых устройств и разных типах протоколов для передачи сообщений между серверной и клиентской частью мобильного приложения. Использование прокси-сервера также позволяет поддержать экзотические нестандартные протоколы, оставшиеся бы непокрытыми стандартными инструментами для сканирования уязвимостей.

